\documentclass[exercice]{thedocument}%
\usepackage{tabularray}%
\usepackage{texmaths-colors}
\nodata
\begin{document}%
    \>Objet1=la droite passant par les points F et G;%
    \>Objet2=[ZX];%
    \>Objet3=%%
    {\begin{tikzpicture}[scale=0.5]
        \coordinate[label=above left:D](A)at(0,0);
        \coordinate[label=above left:E](B)at(3,1);
        \draw node at(A){+};
        \draw node at(B){+};
        \tkzDrawLine[add=0.3 and 0.3](A,B)
    \end{tikzpicture}};%
    \>Objet4=le segment d'extrémités les points H et J;%
    \>Objet5=[PR);%
    \>Objet6=%%
    {\begin{tikzpicture}[scale=0.5]
        \coordinate[label=above left:V](A)at(0,0);
        \coordinate[label=above left:N](B)at(3,1);
        \draw node at(A){+};
        \draw node at(B){+};
        \tkzDrawLine[add=0 and 0.25](A,B)
    \end{tikzpicture}};%
    \enonce{0140}%
\end{document}